% !TEX program = xelatex
\documentclass[a4paper]{article}
\usepackage{amsthm}
\usepackage{amssymb}
\usepackage{bm}
\usepackage{mathtools}
\usepackage[x11names]{xcolor}
\usepackage{xparse}
\usepackage{fontspec}
\usepackage{unicode-math}
\setromanfont{DovesType-Regular.otf}
\setsansfont{Andika}
\setmathfont{Asana Math}[Scale=1]

\usepackage{varwidth}
\usepackage{siunitx}
\usepackage{graphicx}
\usepackage[margin=1cm]{geometry}
\usepackage[most]{tcolorbox}
\usepackage{pgfplots}
\pgfplotsset{compat=newest}
\tcbuselibrary{skins,xparse,poster,breakable}
% \usetikzlibrary{fadings}
\usetikzlibrary{calc, plotmarks, shapes, shapes.geometric, positioning, angles, intersections, quotes, through, patterns, turtle, arrows.meta}
\usetikzlibrary{decorations.markings}
% \usepackage{etoolbox}
\usepackage{tkz-euclide}

%%%%%%%%%%%%%%%%%%%%%%%%%%%%%%%%%%%%%%%%%%%%%%%%%%%%%%%%%
\newcommand\markangle[9]{% origin X Y radius radiusmark mark colour opacity
%  % fill red circle offset-from-centre
  \begin{scope}
    \path[clip] (#1) -- (#2) -- (#3);
    \fill[color=#7,fill opacity=#8,draw=black,name path global=pcircle]  % global declaration required otherwise pcircle is not seen by the `named intersections=' lines below.
    (#1) circle (#4);
  \end{scope}
  % middle calculation
  \path[name path=line one] (#1) -- (#2);
  \path[name path=line two] (#1) -- (#3);
  \path[%
  name intersections={of=line one and pcircle, by={inter one}},
  name intersections={of=line two and pcircle, by={inter two}}
  ] (inter one) -- (inter two) coordinate[pos=#9] (place);
  % put mark
  \node at ($(#1)!#5!(place)$) {\scriptsize{#6}};
}

% \newcommand{\condSoln}[2]{\ifcsdef{r@#1}{#2}{}}

% \newcommand\fadingtext[3][]{%
%    \begin{tikzfadingfrompicture}[name=fading letter]
%      \node[text=transparent!0,inner xsep=0pt,outer xsep=0pt,#1] {#3};
%    \end{tikzfadingfrompicture}%
%    \begin{tikzpicture}[baseline=(textnode.base)]
%      \node[inner sep=0pt,outer sep=0pt,#1](textnode){\phantom{#3}};
%      \shade[path fading=fading letter,#2,fit fading=false]
%      (textnode.south west) rectangle (textnode.north east);%
%    \end{tikzpicture}%
% }

\definecolor{JISpurple}{RGB}{89,72,122}
\definecolor{JISivory}{RGB}{241,234,221}
\definecolor{JIStaupe}{RGB}{183,156,154}
\definecolor{PaleGreen}{RGB}{240,255,240} % 'Honeydew'

\AddToHook{shipout/background}{%
    \put (0in,-\paperheight){\includegraphics[width=\paperwidth,height=\paperheight]{images/R10V5-R.png}}%
}

\newcommand\numberthis{\addtocounter{equation}{1}\tag{\theequation}}

\newtcolorbox{MyOuterBox}{%
  enhanced,
  breakable,
  frame style=JISpurple,
  colback=JISivory,
  colframe=JISpurple,
  title={\includegraphics[width=0.9cm,height=0.9cm]{images/JIS Final Logo FA-02.png}\raisebox{3mm}{\Large{Maths Challenge}\hspace{26em} \Large{\bfseries\sffamily 6}}},
}

\newtcolorbox{MyInnerBox}[2][]{enhanced,%empty,
coltitle=JISpurple,colback=white,
fonttitle=\bfseries\sffamily,
attach boxed title to top left={yshift=-1.5mm},
boxed title style={empty, size=small, top=1mm, bottom=0pt},
varwidth boxed title=0.5\linewidth,
frame code={
  \path (title.east|-frame.north) coordinate (aux);
\path[draw=JISpurple, line width=0.5mm, rounded corners,fill=white]
(frame.west) |- ([xshift=-2.5mm]title.north east) to[out=0, in=180] ([xshift=7.5mm]aux)-|(frame.east)|-(frame.south)-|cycle;
},
title={#2},#1}

\newtcolorbox{MyInnerSplitBox}[2][]{enhanced,%empty,
bicolor,sidebyside,sidebyside align=top seam,
righthand width=5.5cm,colbacklower=white,
coltitle=JISpurple,colback=white,
fonttitle=\bfseries\sffamily,
attach boxed title to top left={yshift=-1.5mm},
boxed title style={empty, size=small, top=1mm, bottom=0pt},
varwidth boxed title=0.5\linewidth,
frame code={
  \path (title.east|-frame.north) coordinate (aux);
\path[draw=JISpurple, line width=0.5mm, rounded corners,fill=white]
(frame.west) |- ([xshift=-2.5mm]title.north east) to[out=0, in=180] ([xshift=7.5mm]aux)-|(frame.east)|-(frame.south)-|cycle;
},
title={#2},#1}


\newtcolorbox{MySolutionBox}{%
  enhanced,
  breakable,
  frame style=JISpurple,
  colback=PaleGreen, colframe=green,
  title={\Large Solution},
  drop fuzzy shadow,
  halign=left,
}

%%%%%%%%%%%%%%%%%%%%%%%%%%%%%%%%%%%%%%%%%%%%%%%%%%
\newtoggle{SOLUTION}
%%% Uncomment the appropriate line below to show solutions %%%
% \toggletrue{SOLUTION}
\togglefalse{SOLUTION}
%%%%%%%%%%%%%%%%%%%%%%%%%%%%%%%%%%%%%%%%%%%%%%%%%


%%%%%%%%%%%%%%%%%%%%%%%%%%%%%%%%%%%%%%%%%%%%%%%%%%
%%%%%%            DOCUMENT BEGINS           %%%%%%
%%%%%%%%%%%%%%%%%%%%%%%%%%%%%%%%%%%%%%%%%%%%%%%%%%
\begin{document}


  \begin{MyOuterBox}
    \iftoggle{SOLUTION}{Here are the full, or partial solutions.
    }{
      Welcome to this week's Maths Challenge!\\
      Everyone is encouraged to take part.\\
      There are 3 to 5 merit points for correct solutions, depending on the difficulty of the problem and how impressed the marker is with your solution.\\
      There are house points for students who submit 7 or more correct solutions in one term.\\
      Solutions must be explained in detail, responses with just the answer may be ignored.\\
      Drop your solution in the box in the staffroom by Tuesday.
    }
       \begin{MyInnerSplitBox}{Year 8 and below}
     Five squares are positioned as shown. The smallest, shaded square has an area of  \(\SI{1}{\square\cm}\). What is the value of \(h\)?
     \iftoggle{SOLUTION}{%conditional output begin
      \begin{MySolutionBox}
        The shaded square has area \(\SI{1}{\square\cm}\) so it has sides of length \(\SI{1}{\cm}\).\par
        Let the length of the side of the largest square be \(a\). Then we can see that the next largest square must have side \(a-1\) because the shaded square on top of it exactly reaches to the top of the largest square. By the same reasoning the third largest square has side \(a-2\) and the smallest unshaded square has side \(a-3\). Now we can find two expressions for the distance across the diagram left to right.
        \begin{align*}
          Y &= X\\
          (a-1) + a &= (a-2) + (a-3) + h\\
          2a - 1 &= 2a -5 + h\\
          -1 &= -5 + h\\
          h &= 4
        \end{align*}
      \end{MySolutionBox}
    }{}%conditional output end
      \tcblower
      \begin{tikzpicture}[scale=0.36]
        \draw (15,0) rectangle (7,8);
        \draw (7,0) rectangle (0,7);
        \draw (0,7) rectangle (6,13);
        \draw (6,13) rectangle (11,8);
        \fill[blue!30] (6.05,7.05) rectangle (6.95,7.95);
        \draw[thin,dotted] (15,8) -- (15,13);
        \draw[<->] (15,9) -- (11,9) node[midway,fill=white,text height=0.5em] {\(h\)};
        \iftoggle{SOLUTION}{%conditional output begin
          \node[left, blue] at (6,7.5) {\(1\)};
          \node[above, blue] at (6.5,8) {\(1\)};
          \node[above, blue] at (11,0) {\(a\)};
          \node[above, blue] at (3.5,0) {\(a-1\)};
          \node[below, blue] at (3,13) {\(a-2\)};
          \node[below, blue] at (8.5,13) {\(a-3\)};
          \draw[<->,blue] (8,0) -- (8,8) node[midway,fill=white] {\(a\)};
          \draw[<->,blue] (5.5,0) -- (5.5,7) node[midway,fill=white] {\(a-1\)};
          \draw[<->,red!70] (0,-0.5) -- (15,-0.5) node[midway,below] {Y};
          \draw[<->,red!70] (0,13.5) -- (15,13.5) node[midway,above] {X};
        }{}%conditional output end
      \end{tikzpicture}
     \end{MyInnerSplitBox}


    \vspace{0.4cm}
          \begin{MyInnerBox}{Year 9 and above}
        The numbers \(4x\), \(2x-3\) and \(4x-3\) are consecutive terms in a linear sequence. Find the value of \(x\).
      \iftoggle{SOLUTION}{%conditional output begin
      \begin{MySolutionBox}
        A linear sequence is like an arithmetic sequence except that the terms of the sequence need not be integers. There is a constant difference \(d\) between terms. If we suppose that \(4x\) is the first term of the sequence we have:\par\vspace{4mm}
          \begin{tabular}{c||c|c|c|c|c}
             n & \(1\) & \(2\) & \(3\) & \(4\) & \(\cdots\)\\ \hline
             \(n^{th}\) term & \(4x\) & \(2x-3\) & \(4x-3\) & & \\
           \end{tabular}\par\vspace{4mm}
          Now, \(d\), is the difference between each successive term, so
          \begin{align*}
            d &= (4x-3)-(2x-3) \qquad\text{  and  }\qquad d = (2x-3) - 4x\\
            \text{So equating,  }\qquad (4x-3)-(2x-3) &= (2x-3) - 4x\\
            2x &= -2x -3\\
            4x &= -3 \implies x= -\frac{3}{4}\\
          \end{align*}
          We have our solution but we should check by substituting in for \(x\).\par\vspace{4mm}
           \begin{tabular}{c||c|c|c|c|c}
             n & \(1\) & \(2\) & \(3\) & \(4\) & \(\cdots\)\\ \hline
             \(n^{th}\)\; term & \(4x\) & \(2x-3\) & \(4x-3\) & & \\
             \(n^{th}\)\; term & \(-3\) & \(-\frac{9}{2}\) & \(-6\) & &\\
           \end{tabular}\par\vspace{4mm}
           We can see that there is a constant difference of \(d=-\frac{3}{2}\).
      \end{MySolutionBox}
    }{}%conditional output end
    \end{MyInnerBox}


  \end{MyOuterBox}

%%%%%%%%%%%%%%%%%%%%%%%%%%%%%%%%%%%%%%%%%%%%%%%%%%
\end{document}



